\section{État qutrit, purification thermochamp double et déficit informationnel}

Pour un lien de poids \(w>0\), soient
\begin{equation}
\label{rm:eq:modular-link-state}
Z(w)=1+\ee^{-w}+\ee^{-2w},
\qquad
\rho(w)=\frac1{Z(w)}\sum_{n=0}^2\ee^{-wn}|n\rangle\langle n|.
\end{equation}
Les poids des deux familles sont
\begin{align}
w_{90}&=90\Dmod
=0.0100077345053979592447310951874\ldots,
\label{rm:eq:modular-w90}\\
w_{120}&=120\Dmod
=0.0133436460071972789929747935832\ldots,
\qquad \frac{w_{120}}{w_{90}}=\frac43.
\label{rm:eq:modular-w120}
\end{align}
L’état fini de bord est
\begin{equation}
\label{rm:eq:modular-product-state}
\rho_A=\rho(w_{90})^{\otimes18}
\otimes\rho(w_{120})^{\otimes18}.
\end{equation}
Son logarithme reproduit \cref{rm:eq:modular-hamiltonian}, avec \(c=18\log Z(w_{90})+18\log Z(w_{120})\).

\begin{theorem}[Purification thermochamp double]
\label{rm:thm:modular-tfd}
Le vecteur
\begin{equation}
\label{rm:eq:modular-tfd-link}
|\operatorname{TFD}(w)\rangle
=\frac1{\sqrt{Z(w)}}\sum_{n=0}^2
\ee^{-wn/2}|n\rangle_A\otimes|n\rangle_{\bar A}
\end{equation}
purifie \(\rho(w)\). Le produit de dix-huit copies pour chacun des poids \(w_{90},w_{120}\) purifie \(\rho_A\) ; dans le système dupliqué, \(S(\rho_A)\) est exactement l’entropie d’intrication.
\end{theorem}

\begin{proof}
En prenant la trace sur le deuxième facteur, l’orthogonalité de \(|n\rangle_{\bar A}\) élimine les termes croisés et laisse
\[
\operatorname{Tr}_{\bar A}
|\operatorname{TFD}(w)\rangle\langle\operatorname{TFD}(w)|
=\rho(w).
\]
La trace partielle commute avec les produits tensoriels, ce qui prouve l’assertion pour les trente-six liens.
\end{proof}

Pour un lien,
\begin{equation}
\label{rm:eq:modular-link-entropy}
S(w)=\log Z(w)
+w\frac{\ee^{-w}+2\ee^{-2w}}{Z(w)}.
\end{equation}
En conséquence,
\begin{equation}
\label{rm:eq:modular-entropy-value}
S(\rho_A)=18S(w_{90})+18S(w_{120})
=39.54837320881807994957319730\ldots\ \text{nats}.
\end{equation}
L’état maximalement mélangé de trente-six qutrits possède l’entropie \(36\log3\). Définissons le défaut orienté
\begin{equation}
\label{rm:eq:modular-oriented-information}
\boxed{
I_{\rm or}=36\log3-S(\rho_A)
=0.001669183233868940655631225913\ldots\ \text{nats}.}
\end{equation}
Sa positivité découle de la maximalité entropique de l’état uniforme. La proximité de \(I_{\rm or}\) et de \(\gammaH a_0\) ne constitue pas une identité et n’est pas utilisée pour sélectionner l’une ou l’autre des deux valeurs.
